\section{Discussion}

\paragraph{Axis A: Temporal ordering drives the representation gap.}
The C-index difference between Setting 4 (decoupled age--event encoding,
0.6233) and Setting 3 (prose serialisation, 0.6156), using identical
Qwen3-8B encoders and raw features, is the cleanest signal in our
results.
Prose serialisation collapses temporal ordering into semantic proximity:
the sentence \textit{``At age 44, J45''} and \textit{``At age 52, F32''}
differ only in surface tokens, but the 8-year interval between them is
invisible to the encoder's attention mechanism.
Decoupled encoding preserves this structure explicitly: the 128-dim
sinusoidal age vector ensures that events at age 44 and age 52 are
geometrically separated before the token embedding is appended.
This aligns with the structured multimodal intelligence thesis:
\emph{surface semantic encoding is insufficient when temporal causal
structure drives the outcome.}
The +0.008 C-index advantage of S4 over S3 is the primary empirical claim;
bootstrap confidence intervals confirm this gap exceeds sampling noise:
CI\,=\,($+$0.005, $+$0.010), $p<0.001$.
Setting 2 (Delphi's ordinal token stream) falls between S3 and S4:
temporal ordering is preserved as positional structure in the input
sequence, but the autoregressive hazard chain cannot separate
ordering effects from pre-training priors (see calibration gap below).

\paragraph{Setting 5: LLMs encode numeric biomarker structure from prose.}
Setting 5 (unified prose, no raw features) achieves C\,=\,0.6195 ---
matching the 39-feature CoxPH baseline (0.6191) while using no raw
numeric features.
This suggests that Qwen3-8B's pre-training has internalised enough
biomarker-range semantics to distinguish, e.g., HbA1c\,=\,52 from
HbA1c\,=\,38 in terms of mortality relevance.
The near-term TD-AUC advantage (0.605 vs.\ 0.584 at 9.5\,yr) further
suggests the embedding captures interactions between biomarker values and
disease history that a linear CoxPH on raw features misses.
However, this encoding is \emph{implicit} and cannot be inspected or
enforced, which limits trustworthiness in high-stakes clinical deployment.

\paragraph{Ablation axes B and C.}
The encoder and fusion ablations confirm that the Axis~A findings are robust
rather than artefacts of a specific implementation choice.
For Axis~B, BCB and Qwen3-8B produce near-identical C-indices for S3 and S4,
confirming that discriminative information is accessible to both encoder
families; the temporal-ordering advantage of S4 over S3 persists
($\Delta$C\,=\,$+$0.008 for Qwen, $p<0.001$).
The anomalously high S5·BCB C-index (0.6371) is explained by 512-token
truncation retaining primarily age signal, which collapses TD-AUC to 0.507.
For Axis~C, concatenation outperforms summation by $+$0.028--$+$0.039 C-index
across all cells, confirming that retaining independent embedding dimensions
alongside raw features preserves more representational capacity than projective
fusion.
The primary conclusion — that serialisation format drives discrimination more
than encoder family or fusion method — is supported across both ablation axes.

\paragraph{Calibration gap of autoregressive generation.}
Delphi's IBS of 0.185 vs.\ $\approx$0.141 for supervised CoxPH variants
must be interpreted alongside an important asymmetry: Delphi operates solely
on serialised ICD codes and demographic tokens and has no access to the
numeric biomarker features (HbA1c, CRP, creatinine, cystatin C) available to
the CoxPH baselines.
This input asymmetry means the gap partly reflects missing acute biomarker
signal rather than a deficiency of the autoregressive approach within its
intended operating regime.
More fundamentally, the sigmoid hazard chain produces survival curves without
a calibration reference from the target population, an inherent constraint of
zero-shot deployment under distribution shift between pre-training and the
UK Biobank cohort.
The near-term AUC collapse (0.52 at 9.5\,yr) is consistent with this
interpretation: short-horizon mortality risk is dominated by derangements in
CRP, creatinine, and albumin that do not appear in the ICD token stream.
These results therefore reflect a structural limitation of the autoregressive
generation paradigm for \emph{acute} risk stratification: the paradigm
lacks access to the quantitative physiological signal that most strongly
predicts near-term death.
Fine-tuning with a calibration-aware objective and numeric biomarker inputs
would likely close this gap substantially, but at the cost of zero-shot
generalisability across cohorts.

\paragraph{Limitations.}
All embedding methods freeze their encoders; fine-tuning may yield
additional gains at substantially higher compute cost.
IBS and TD-AUC rely on independent censoring assumptions that may be
violated in UKB.
All results are from a single biobank cohort; external validation is
required before clinical translation.
Bio\_ClinicalBERT's 512-token limit is an architectural constraint that
disproportionately affects Setting 5; future work should evaluate
longer-context clinical encoders.
The primary evaluation metrics (C-index, TD-AUC) measure discrimination but not
calibration of the predicted survival function.
We supplement IBS with restricted mean survival time (RMST\,$= \int_0^\tau S(t\mid x)\,\mathrm{d}t$),
which provides a patient-interpretable calibration summary (expected life-years
within the follow-up window $\tau$).
RMST-decile calibration analysis across the full 12-cell matrix is left as
an extension once all server embedding runs are consolidated locally.
